\documentclass[11pt]{article}
\usepackage{amsmath,amssymb}

\setlength{\parindent}{0pt}
\setlength{\parskip}{0.7em}
\setlength{\textwidth}{6.5in}
\setlength{\oddsidemargin}{0in}
\setlength{\evensidemargin}{0in}
\setlength{\topmargin}{-0.4in}
\setlength{\textheight}{9.2in}

\begin{document}

\begin{center}
{\Large SYSC 5108 Exam Study: Assignment 3, Question 3}\\
\vspace{0.3em}
{\large Forward propagation through a small neural network}\\
\vspace{0.3em}
Source question: \texttt{SYSC 5108 W26 Assignment 3.pdf}, Question 3
\end{center}

\section*{Question Summary}

The assignment gives a feedforward neural network with two sensing/input neurons and three processing neurons:
\[
\text{inputs: neurons }1,2,\qquad
\text{processing neurons: }3,4,5.
\]
The input is
\[
x_1=0.7,\qquad x_2=0.3.
\]
From the network figure and the submitted solution, the weights and biases are:
\[
\begin{array}{c|ccc}
\text{Neuron} & \text{bias} & \text{weight from }x_1\text{ or }a_3 & \text{weight from }x_2\text{ or }a_4\\
\hline
3 & w_{3,0}=0.6 & w_{3,1}=0.1 & w_{3,2}=0.5\\
4 & w_{4,0}=0.8 & w_{4,1}=0.2 & w_{4,2}=0.4\\
5 & w_{5,0}=0.9 & w_{5,3}=0.3 & w_{5,4}=0.7
\end{array}
\]

The question asks for:
\begin{itemize}
\item the output of neuron 5 if all processing neurons use logistic activation,
\item the output of neuron 5 if all processing neurons use ReLU activation,
\item the matrix form for the ReLU case.
\end{itemize}

\section*{Step 1: Recall the Forward Propagation Rule}

For each processing neuron, first compute the weighted sum:
\[
z_j=w_{j,0}+\sum_i w_{j,i}a_i.
\]
Then apply the activation function:
\[
a_j=\phi(z_j).
\]

Chapter 4 slides on deep feedforward networks use the same idea: each layer applies a linear transformation followed by a nonlinear activation. Goodfellow, Bengio, and Courville describe a feedforward network as a composition of layers; for a hidden layer, the standard form is a matrix multiplication plus bias followed by an activation function.

The two activation functions needed here are:
\[
\sigma(z)=\frac{1}{1+e^{-z}}
\]
for logistic activation, and
\[
\operatorname{ReLU}(z)=\max(0,z)
\]
for rectified linear activation.

\section*{Step 2: Compute Hidden-Layer Weighted Sums}

Neuron 3:
\[
z_3=w_{3,0}+w_{3,1}x_1+w_{3,2}x_2.
\]
Substitute:
\[
z_3=0.6+(0.1)(0.7)+(0.5)(0.3).
\]
So
\[
z_3=0.6+0.07+0.15=0.82.
\]

Neuron 4:
\[
z_4=w_{4,0}+w_{4,1}x_1+w_{4,2}x_2.
\]
Substitute:
\[
z_4=0.8+(0.2)(0.7)+(0.4)(0.3).
\]
So
\[
z_4=0.8+0.14+0.12=1.06.
\]

\section*{Part (a): Logistic Activation}

With logistic activation:
\[
a_3=\sigma(z_3)=\sigma(0.82)
=\frac{1}{1+e^{-0.82}}
\approx 0.694236.
\]
Similarly:
\[
a_4=\sigma(z_4)=\sigma(1.06)
=\frac{1}{1+e^{-1.06}}
\approx 0.742691.
\]

Now compute the weighted sum for neuron 5:
\[
z_5=w_{5,0}+w_{5,3}a_3+w_{5,4}a_4.
\]
Substitute:
\[
z_5=0.9+(0.3)(0.694236)+(0.7)(0.742691).
\]
So
\[
z_5\approx 0.9+0.208271+0.519884=1.628154.
\]
Apply logistic activation at neuron 5:
\[
a_5=\sigma(z_5)=\sigma(1.628154)
\approx 0.835917.
\]
Therefore:
\[
\boxed{a_5\approx 0.836\quad\text{for logistic activation}.}
\]

\section*{Part (b): ReLU Activation}

With ReLU activation:
\[
a_3=\operatorname{ReLU}(0.82)=0.82,
\qquad
a_4=\operatorname{ReLU}(1.06)=1.06.
\]
Both values are unchanged because both weighted sums are positive.

Now compute neuron 5:
\[
z_5=0.9+(0.3)(0.82)+(0.7)(1.06).
\]
So
\[
z_5=0.9+0.246+0.742=1.888.
\]
Apply ReLU:
\[
a_5=\operatorname{ReLU}(1.888)=1.888.
\]
Therefore:
\[
\boxed{a_5=1.888\quad\text{for ReLU activation}.}
\]

\section*{Part (c): Matrix Form for the ReLU Network}

Write the input vector as
\[
x=
\begin{bmatrix}
0.7\\
0.3
\end{bmatrix}.
\]
For the hidden layer containing neurons 3 and 4:
\[
W^{(1)}=
\begin{bmatrix}
0.1 & 0.5\\
0.2 & 0.4
\end{bmatrix},
\qquad
b^{(1)}=
\begin{bmatrix}
0.6\\
0.8
\end{bmatrix}.
\]
Compute the hidden pre-activation vector:
\[
z^{(1)}=W^{(1)}x+b^{(1)}.
\]
Substitute:
\[
z^{(1)}
=
\begin{bmatrix}
0.1 & 0.5\\
0.2 & 0.4
\end{bmatrix}
\begin{bmatrix}
0.7\\
0.3
\end{bmatrix}
+
\begin{bmatrix}
0.6\\
0.8
\end{bmatrix}
=
\begin{bmatrix}
0.82\\
1.06
\end{bmatrix}.
\]
Apply ReLU elementwise:
\[
a^{(1)}=\operatorname{ReLU}(z^{(1)})
=
\begin{bmatrix}
0.82\\
1.06
\end{bmatrix}.
\]

For the output layer containing neuron 5:
\[
W^{(2)}=
\begin{bmatrix}
0.3 & 0.7
\end{bmatrix},
\qquad
b^{(2)}=0.9.
\]
Compute:
\[
z^{(2)}=W^{(2)}a^{(1)}+b^{(2)}.
\]
Substitute:
\[
z^{(2)}
=
\begin{bmatrix}
0.3 & 0.7
\end{bmatrix}
\begin{bmatrix}
0.82\\
1.06
\end{bmatrix}
+0.9
=1.888.
\]
Apply ReLU:
\[
a^{(2)}=\operatorname{ReLU}(1.888)=1.888.
\]

\section*{Step 3: Why the Logistic and ReLU Answers Differ}

The logistic activation squashes all neuron outputs into the interval $(0,1)$. Therefore the final output in part (a) is also between 0 and 1:
\[
a_5\approx 0.836.
\]
ReLU does not squash positive values. It keeps positive weighted sums unchanged and only clips negative values to zero. Since every pre-activation here is positive, the ReLU network behaves like a sequence of affine transformations, giving the larger output:
\[
a_5=1.888.
\]

\section*{Concise Exam Answer}

\[
z_3=0.6+0.1(0.7)+0.5(0.3)=0.82,\qquad
z_4=0.8+0.2(0.7)+0.4(0.3)=1.06.
\]
For logistic activation:
\[
a_3=\sigma(0.82)=0.6942,\qquad
a_4=\sigma(1.06)=0.7427,
\]
\[
z_5=0.9+0.3(0.6942)+0.7(0.7427)=1.6282,
\qquad
a_5=\sigma(1.6282)=0.8359.
\]
For ReLU activation:
\[
a_3=0.82,\qquad a_4=1.06,
\]
\[
z_5=0.9+0.3(0.82)+0.7(1.06)=1.888,
\qquad
a_5=1.888.
\]
The ReLU matrix form is:
\[
z^{(1)}=
\begin{bmatrix}
0.1 & 0.5\\
0.2 & 0.4
\end{bmatrix}
\begin{bmatrix}
0.7\\
0.3
\end{bmatrix}
+
\begin{bmatrix}
0.6\\
0.8
\end{bmatrix}
=
\begin{bmatrix}
0.82\\
1.06
\end{bmatrix},
\]
\[
a^{(1)}=\operatorname{ReLU}(z^{(1)})
=
\begin{bmatrix}
0.82\\
1.06
\end{bmatrix},
\]
\[
z^{(2)}=
\begin{bmatrix}
0.3 & 0.7
\end{bmatrix}
a^{(1)}+0.9
=1.888,
\qquad
a^{(2)}=1.888.
\]

\section*{Cross References Used}

\begin{itemize}
\item Assignment: \texttt{SYSC 5108 W26 Assignment 3.pdf}, Question 3.
\item Local submitted Assignment 3 PDF checked for comparison, Question 3 response.
\item Slides: Chapter 4 slide deck on deep feedforward networks, hidden layers, nonlinear activations, ReLU, and logistic sigmoid units.
\item Textbook: Goodfellow, Bengio, and Courville, \textit{Deep Learning}, Chapter 6, Section 6.1 on deep feedforward networks.
\item Textbook: Goodfellow, Bengio, and Courville, \textit{Deep Learning}, Chapter 6, discussion around feedforward network layers, matrix multiplication, activations, and ReLU units.
\item Textbook: Goodfellow, Bengio, and Courville, \textit{Deep Learning}, Chapter 3, Section 3.10 on the logistic sigmoid.
\end{itemize}

\end{document}
